% TOP_PROBLEM: {"id": 4700002, "title": "Systems with homogeneous components I", "category_id": 11, "category": {"id": 11, "name": "dynamical_systems", "display_name": "Dynamical Systems", "description": "Problems about long-term behavior of deterministic systems, Hamiltonian dynamics, and periodic orbits.", "slug": "dynamical-systems", "order_index": 11, "created_at": "2026-07-31T15:26:25.671Z"}, "status": "open", "problem_number": "AMR-046-0002", "set_id": 13}
% FIRST_POSTED: 2026-10-02
% ENTRY_KIND: statement_only
% UnsolvedMath ID 4700002; source catalogue code AMR-046-0002
% !TeX program = lualatex
% Definition and sources only; this file is not an LLM solution attempt.
\documentclass[11pt,a4paper]{article}
\usepackage[margin=25mm]{geometry}
\usepackage{fontspec}
\setmainfont{lmroman10-regular.otf}[BoldFont=lmroman10-bold.otf,
 ItalicFont=lmroman10-italic.otf,BoldItalicFont=lmroman10-bolditalic.otf]
\usepackage{amsmath,amssymb,mathtools}
\usepackage{unicode-math}
\setmathfont{latinmodern-math.otf}
\AtBeginDocument{\let\setminus\smallsetminus}
\usepackage{xurl,hyperref}
\hypersetup{colorlinks=true,urlcolor=blue}
\setcounter{secnumdepth}{0}
\setlength{\parindent}{0pt}
\setlength{\parskip}{5pt}
\setlength{\emergencystretch}{3em}
\begin{document}
\section*{Systems with homogeneous components I}\label{4700002}
{\small UnsolvedMath ID 4700002 · AMR-046-0002 · Dynamical Systems · AMR Open Problem Lists}
\subsection{Definitions and mathematical statement}
Fix distinct positive odd integers $n,m$. A real polynomial $P_n$ is homogeneous
of degree $n$ when $P_n(tx,ty)=t^nP_n(x,y)$ for every real $t$; require $P_n$
and $Q_m$ to be nonzero and of the indicated degrees. Consider
\[
 \dot x=P_n(x,y),\qquad \dot y=Q_m(x,y).
\]
A periodic orbit is the closed image of a nonconstant solution. It is a
limit cycle if a neighborhood of that orbit contains no other periodic orbit.
Let $L(n,m)$ be the supremum, over the coefficients of both homogeneous
polynomials, of the number of distinct limit cycles in the whole plane.
The question is whether
\[
 L(n,m)=\frac{n+m}{2}\qquad(n,m\text{ odd},\ n\ne m).
\]
In particular, the required upper bound must be independent of the sizes and
signs of the coefficients and of the locations of the cycles. Cycles are
counted geometrically, without their algebraic or bifurcation multiplicities.

The restriction to unequal odd degrees is part of the problem. Equal-degree
homogeneity rescales periodic orbits into continuous families, and those
families do not consist of isolated cycles. The source establishes examples
with at least $(n+m)/2$ cycles when the degrees are unequal and odd. Hence
the remaining assertion is that this construction already attains the global
maximum. A proof for a neighborhood of the origin or a particular perturbation
family would not establish the stated bound for all coefficients.
\subsection{Short English statement}
Is the known lower bound of half the sum of the two odd degrees also the exact global maximum number of limit cycles?
\subsection{Sources}
\begin{itemize}
\item[S1] ULAM.AI and the UnsolvedMath Contributors, supplied problems.json, record 4700002 (AMR-046-0002), AMR Open Problem Lists. Catalogue identity and source transcription; CC BY 4.0.
\url{https://huggingface.co/datasets/ulamai/UnsolvedMath}.
\item[S2] Gasull - Some open problems in low dimensional dynamical systems (2020), Problem environment 2: Systems with homogeneous components I. Original source indexed by AMR-046-0002.
\url{https://arxiv.org/abs/2012.02524}.
\end{itemize}
\end{document}
